\documentclass[11pt,a4paper]{article}

% ── Codificación y fuentes ────────────────────────────────────────────────────
\usepackage[utf8]{inputenc}
\usepackage[T1]{fontenc}
\usepackage{lmodern}
\usepackage[spanish]{babel}

% ── Geometría y espaciado ─────────────────────────────────────────────────────
\usepackage[top=2.5cm, bottom=2.5cm, left=2.8cm, right=2.8cm]{geometry}
\usepackage{setspace}
\setstretch{1.15}
\usepackage{parskip}

% ── Tablas ────────────────────────────────────────────────────────────────────
\usepackage{booktabs}
\usepackage{tabularx}
\usepackage{array}
\usepackage{longtable}
\usepackage{enumitem}

% ── Colores y cajas ───────────────────────────────────────────────────────────
\usepackage[dvipsnames]{xcolor}
\usepackage{tcolorbox}
\tcbuselibrary{skins, breakable}

% ── Cabeceras y pies de página ────────────────────────────────────────────────
\usepackage{fancyhdr}
\usepackage{lastpage}
\pagestyle{fancy}
\fancyhf{}
\fancyhead[L]{\small\textbf{Informe de Evaluación} --- Física para Bioanalistas}
\fancyhead[R]{\small karla rodriguez 33628014}
\fancyfoot[C]{\small Página \thepage\ de \pageref{LastPage}}
\renewcommand{\headrulewidth}{0.4pt}

% ── Hiperenlaces ──────────────────────────────────────────────────────────────
\usepackage[colorlinks=true, linkcolor=NavyBlue, urlcolor=NavyBlue]{hyperref}

% ── Paleta de colores personalizados ─────────────────────────────────────────
\definecolor{desempeno_excelente}{HTML}{1A6B2F}
\definecolor{desempeno_bueno}{HTML}{2E5A9C}
\definecolor{desempeno_suficiente}{HTML}{8B6914}
\definecolor{desempeno_deficiente}{HTML}{C0392B}
\definecolor{desempeno_insuficiente}{HTML}{7B241C}
\definecolor{grisFondo}{HTML}{F5F5F5}
\definecolor{azulTitulo}{HTML}{1B2A6B}
\definecolor{verdeFortaleza}{HTML}{1A6B2F}
\definecolor{naranjaMejora}{HTML}{A04000}

% ── Estilos de cajas ─────────────────────────────────────────────────────────
\tcbset{
  cajaTitulo/.style={
    enhanced, breakable,
    colback=azulTitulo!8, colframe=azulTitulo,
    fonttitle=\bfseries\large, coltitle=white,
    attach boxed title to top left={yshift=-2mm, xshift=4mm},
    boxed title style={colback=azulTitulo},
    top=6pt, bottom=6pt, left=8pt, right=8pt,
  },
  cajaFortaleza/.style={
    enhanced, breakable,
    colback=verdeFortaleza!6, colframe=verdeFortaleza!60,
    fonttitle=\bfseries, coltitle=verdeFortaleza,
    top=4pt, bottom=4pt, left=8pt, right=8pt,
  },
  cajaMejora/.style={
    enhanced, breakable,
    colback=naranjaMejora!6, colframe=naranjaMejora!60,
    fonttitle=\bfseries, coltitle=naranjaMejora,
    top=4pt, bottom=4pt, left=8pt, right=8pt,
  },
  cajaRecomendacion/.style={
    enhanced, breakable,
    colback=grisFondo, colframe=gray!50,
    fonttitle=\bfseries, coltitle=black,
    top=4pt, bottom=4pt, left=8pt, right=8pt,
  },
}

% ─────────────────────────────────────────────────────────────────────────────
\begin{document}

% ══ PORTADA ══════════════════════════════════════════════════════════════════
\begin{center}
  {\Large\bfseries\color{azulTitulo} Universidad de Oriente/Núcleo Sucre --- Física para Bioanalistas}\\[4pt]
  {\large Informe de Evaluación Preliminar}\\[12pt]
  \rule{\linewidth}{1.2pt}\\[6pt]
  {\LARGE\bfseries karla rodriguez 33628014}\\[4pt]
  \rule{\linewidth}{0.6pt}
\end{center}

% ══ FICHA TÉCNICA ════════════════════════════════════════════════════════════
\vspace{4pt}
\begin{tcolorbox}[cajaTitulo, title=Datos del informe]
\begin{tabular}{@{}llll@{}}
  \textbf{Estudiante:}  & karla rodriguez 33628014  &
  \textbf{Fecha:}       & 2026-06-26 \\[4pt]
  \textbf{Archivo PDF:} & \texttt{karla\_rodriguez\_33628014.pdf}  &
  \textbf{Páginas:}     & 5 \\
\end{tabular}
\end{tcolorbox}

% ══ RESULTADO GLOBAL ═════════════════════════════════════════════════════════
\vspace{6pt}
\begin{center}
  \begin{tcolorbox}[
    enhanced,
    colback=white, colframe=desempeno_bueno,
    width=0.62\linewidth,
    boxrule=2pt,
    halign=center, valign=center,
    top=8pt, bottom=8pt,
  ]
    \centering
    {\large\bfseries Puntaje sugerido}\\[6pt]
    {\Huge\bfseries\color{desempeno_bueno} 8.5/10}\\[6pt]
    {\large Nivel de desempeño: \textbf{\color{desempeno_bueno} Bueno}}
  \end{tcolorbox}
\end{center}

% ══ RESUMEN GENERAL ══════════════════════════════════════════════════════════
\section*{Resumen general}
El estudiante mostró un buen dominio de los conceptos físicos y procedimientos matemáticos, aunque cometió algunos errores procedimentales en cálculos. La claridad y organización de las respuestas fueron adecuadas.

% ══ FORTALEZAS Y ÁREAS DE MEJORA ════════════════════════════════════════════
\vspace{4pt}
\begin{minipage}[t]{0.48\linewidth}
  \begin{tcolorbox}[cajaFortaleza, title={Fortalezas}]
\begin{itemize}[leftmargin=1.5em, itemsep=2pt]
  \item Correcta aplicación del principio de Arquímedes y cálculo de volumen de muestra
  \item Uso adecuado de la ecuación de Bernoulli para determinar presión en arteria estrechada
  \item Comprensión del concepto de ley de Poiseuille y cálculo de caudal disminuido
  \item Identificación correcta de la relación entre viscosidad y resistencia al flujo
\end{itemize}
  \end{tcolorbox}
\end{minipage}
\hfill
\begin{minipage}[t]{0.48\linewidth}
  \begin{tcolorbox}[cajaMejora, title={Areas de mejora}]
\begin{itemize}[leftmargin=1.5em, itemsep=2pt]
  \item Errores en el cálculo de presión hidrostática (confusión entre presión y altura)
  \item Cálculo incorrecto de presión en arteria estrechada (error en aplicación de Bernoulli)
  \item Uso inconsistente de unidades en algunos cálculos
  \item Falta de redondeo adecuado en resultados finales
\end{itemize}
  \end{tcolorbox}
\end{minipage}

% ══ OBSERVACIONES POR PREGUNTA ═══════════════════════════════════════════════
\section*{Observaciones por pregunta}

\begin{longtable}{>{\bfseries}p{2.8cm} p{1.6cm} p{7.0cm} p{2.8cm}}
  \toprule
  \textbf{Pregunta} & \textbf{Puntaje} & \textbf{Evaluación} & \textbf{Errores detectados} \\
  \midrule
  \endhead
  \bottomrule
  \endfoot
    \textbf{1. Presión Hidrostática y Venosa} & 1.5/2.0 & El estudiante resolvió correctamente la ecuación de presión hidrostática, pero cometió un error en la aplicación de la fórmula (confundió presión con altura). & \textit{Confusión entre presión y altura en la fórmula h = Pv / ($\rho$g)} \\
    \midrule
    \textbf{2. Principio de Arquímedes} & 2.0/2.0 & Respuesta correcta y completa, con cálculos precisos y unidades adecuadas. & \textit{---} \\
    \midrule
    \textbf{3. Hemodinámica (Continuidad y Bernoulli)} & 1.5/2.0 & El estudiante aplicó correctamente la ecuación de continuidad, pero cometió un error en la aplicación de la ecuación de Bernoulli para calcular la presión en la arteria estrechada. & \textit{Aplicación incorrecta de la ecuación de Bernoulli (error en cálculo de presión estática)} \\
    \midrule
    \textbf{4. Ley de Poiseuille} & 2.0/2.0 & Respuesta correcta y completa, con cálculos precisos y unidades adecuadas. & \textit{---} \\
    \midrule
    \textbf{5. Flujo Viscoso y Resistencia} & 1.5/2.0 & El estudiante aplicó correctamente la fórmula de resistencia al flujo, pero cometió un error en el cálculo de la diferencia de presión ($\Delta$P). & \textit{Cálculo incorrecto de la diferencia de presión ($\Delta$P)} \\
    \midrule
\end{longtable}

% ══ ERRORES DETECTADOS ═══════════════════════════════════════════════════════
\section*{Análisis de errores}

\subsection*{Errores conceptuales}
\begin{itemize}[leftmargin=1.5em, itemsep=2pt]
  \item Confusión entre presión y altura en la fórmula de presión hidrostática
  \item Aplicación incorrecta de la ecuación de Bernoulli para calcular presión en arteria estrechada
\end{itemize}

\subsection*{Errores procedimentales}
\begin{itemize}[leftmargin=1.5em, itemsep=2pt]
  \item Error en cálculo de presión en arteria estrechada (fórmula de Bernoulli)
  \item Cálculo incorrecto de la diferencia de presión en flujo viscoso
  \item Uso inconsistente de unidades en algunos cálculos
\end{itemize}

% ══ RECOMENDACIONES AL ESTUDIANTE ════════════════════════════════════════════
\section*{Retroalimentación para el estudiante}
\begin{tcolorbox}[cajaRecomendacion, title={Dirigido a: karla rodriguez 33628014}]
Revisar cuidadosamente la aplicación de las fórmulas físicas, especialmente en problemas de presión y flujo. Prestar atención a la consistencia en el uso de unidades y al redondeo adecuado de resultados finales.
\end{tcolorbox}

% ══ NOTA PARA EL PROFESOR ════════════════════════════════════════════════════
\section*{Nota para el profesor}
\begin{tcolorbox}[
  enhanced, breakable,
  colback=yellow!8, colframe=yellow!60!black,
  fonttitle=\bfseries, coltitle=black,
  title={Observaciones para revisión manual},
  top=4pt, bottom=4pt, left=8pt, right=8pt,
]
El examen requiere revisión cuidadosa, especialmente en las preguntas 1 y 3 donde se identificaron errores conceptuales y procedimentales. Las respuestas a las preguntas 2, 4 y 5 fueron correctas y completas.
\end{tcolorbox}

% ══ PIE DE INFORME ════════════════════════════════════════════════════════════
\vspace{12pt}
\begin{center}
  \footnotesize\color{gray}
  Informe generado automáticamente por el sistema de evaluación con IA (nvidia/nemotron-nano-12b-v2-vl:free) y soporte LaTex. \\
  Este documento es una evaluación \textbf{preliminar} para ser revisado por el profesor antes de ser entregado al 
  estudiante. \\
  Generado el 2026-06-26 \textbullet\ BIOANALISIS Sistema Académico de Evaluación.
  \end{center}

\end{document}
